\textbf{Determine planet periods.}
Synodic period of planet 1 and 2 (assume $P_1 < P_2$) is
\[ \frac{1}{P_S} = \frac{1}{P_1} - \frac{1}{P_2}
   \quad\Rightarrow\quad
   P_S = \frac{P_1 P_2}{P_2 - P_1} \]
Resonant orbit happens if
\[ P_S = mP_1
   \quad\Rightarrow\quad
   \frac{P_1 P_2}{P_2 - P_1} = mP_1
   \quad\Rightarrow\quad
   m = \frac{P_2}{P_2 - P_1} > 1,\ \text{integer} \]

Using the data of the Gliesean system we have
\hfill(50 points)

\begin{center}
\begin{tabular}{|l|c|c|c|}
\hline
\textbf{Gliese System} & \textbf{Mass (kg)} & \textbf{Semi Major Axis (m)}
  & \textbf{Period (Earth, days)} $P = 2\pi\sqrt{\dfrac{a^3}{GM_G}}$ \\
\hline
Gliese 876 star & $6.64359 \times 10^{29}$ & & \\
\hline
Gliese 876 b & $4.31938 \times 10^{27}$ & $3.1161 \times 10^{10}$ & 60.07568 \\
\hline
Gliese 876 c & $1.35564 \times 10^{27}$ & $1.9388 \times 10^{10}$ & 29.48304 \\
\hline
Gliese 876 d & $4.079 \times 10^{25}$ & $3.1116 \times 10^{9}$ & 1.8956 \\
\hline
Gliese 876 e & $8.7194 \times 10^{25}$ & $5.0011 \times 10^{10}$ & 122.1432 \\
\hline
\end{tabular}
\end{center}

\textbf{Calculate resonant orbit for each planet.}
From the results above we see that
\begin{itemize}
\item Gliese 876 c and Gliese 876 b have resonant orbit
  \[ m = \frac{P_b}{P_b - P_c} = 1.96 \approx 2 \]
  \hfill(25 points)
\item Gliese 876 b and Gliese 876 e have resonant orbit
  \[ m = \frac{P_e}{P_e - P_b} = 1.9679 \approx 2 \]
  \hfill(25 points)
\end{itemize}
